\chapter{Command-line reference}
\label{ch:cli}

\section{Synopsis}

\begin{codeblock}
fblockkit                      interactive menu (Chapter~\ref{ch:running})
fblockkit --record FILE        the same, recording this session's input to FILE
fblockkit run SCRIPT.txt       replay a script
fblockkit search KEYWORDS      search the tool index (Chapter~\ref{ch:menu5})
fblockkit guide TOOL_ID        tool onboarding notes (Chapter~\ref{ch:menu6})
\end{codeblock}

All output goes to standard output; there is no log file and no state between
runs. Exit codes: \code{0} success, \code{2} usage problem (missing script,
empty keyword, unknown tool id). The interactive menu is also reachable as
\code{python -m fblockkit} when the package is used from source.

\section{The script format}

One input per line. An empty line means ``press Enter'' (accept the default);
a line starting with \code{\#} is a comment and is skipped; everything else is
delivered to the prompt exactly as typed. Scripts can be produced by
\code{--record}, by \menu{8}, or by hand.

Replaying is exactly equivalent to typing: same prompts, same output,
byte-identical across runs. A script that runs out of lines mid-session ends
with ``Cancelled (input finished)'' for the pending prompt and continues to the
next menu (or exits cleanly at EOF).

\section{Examples}

\begin{codeblock}
fblockkit search lanthanide DMET
fblockkit guide openmolcas
fblockkit run my_session.txt
fblockkit --record my_session.txt      # then work normally; the file is written at exit
\end{codeblock}

\section{Environment}

\code{PYTHONUTF8=1} is recommended on Windows consoles (the program prints
ASCII/UTF-8 only, but the console code page can mangle the display). On the
supported Python versions the program has no other environment requirements.
